% GENERATED FILE. Edit canonical lessons, then run npm run generate:books.
% canonical-source-hash: fnv1a64:4bd2ccbd60a194e7
% canonical-lessons: HI-C423-gard, HI-C423-khambha, HI-C423-sor, HI-C423-gurrana, HI-C423-bhaukna

\chapter{Guards and Street Dogs}
\label{ch:hi-guards-and-street-dogs}
\hlchaptermodality{423}

\begin{chapteropening}
\textbf{By the end of this chapter:} \emph{I can talk about the guard and the lamp posts, the noise of the street and the dogs at night.}

Complete the last lesson of chapter 423: I can talk about the guard and the lamp posts, the noise of the street and the dogs at night.
\end{chapteropening}

\section[\texorpdfstring{\hi{गार्ड} (gārḍ)}{गार्ड (gārḍ)}]{\hi{गार्ड} (gārḍ) — a guard}

\label{lesson:HI-C423-gard}



\begin{quote}
Before the new one: say the Hindi for a compound, then the Hindi for a thatched roof.
\end{quote}

\subsection*{You'll want to know: \hi{गार्ड}}
\textbf{\hi{गार्ड}} — \emph{gārḍ} — \textquotedblleft{}a guard\textquotedblright{}.

\begin{quote}
\hi{बैंक} \hi{के} \hi{दरवाज़े} \hi{पर} \hi{एक} \hi{गार्ड} \hi{खड़ा} \hi{था।}
\end{quote}

\begin{quote}
\emph{baiṅk ke darvāze par ek gārḍ khaṛā thā.}
\end{quote}

\begin{quote}
A guard was standing at the door of the bank.
\end{quote}

Masculine or feminine, by the person meant.

\begin{cousinweb}
\textbf{\hi{गार्ड}} (\emph{gārḍ}) is borrowed from English \emph{guard}, the security guard at a bank, a shop or a block of flats; on the railway it is also the guard at the back of the train.
\end{cousinweb}

\begin{grammarlens}[title={What you've built}]
One more word for getting around.
\end{grammarlens}

\subsection*{Your turn}
\begin{itemize}
\raggedright
  \item \emph{Say it:} \emph{gārḍ}
  \item \emph{Say it:} \emph{gārḍ}, once more
  \item \emph{Say it:} \emph{baiṅk ke darvāze par ek gārḍ khaṛā thā.} — A guard was standing at the door of the bank.
  \item \emph{Recall:} say the Hindi for a compound, then the Hindi for a thatched roof, then say \emph{gārḍ} again
\end{itemize}

\begin{tcolorbox}[breakable,colback=teal!4,colframe=teal!35!black,title={Before you move on}]
What does \hi{गार्ड} mean? (\textquotedblleft{}A guard\textquotedblright{}.) Say it once more.
\end{tcolorbox}
\section[\texorpdfstring{\hi{खंभा} (khambhā)}{खंभा (khambhā)}]{\hi{खंभा} (khambhā) — a pole; a pillar}

\label{lesson:HI-C423-khambha}



\begin{quote}
Before the new one: say the Hindi for a thatched roof, then the Hindi for a guard.
\end{quote}

\subsection*{You'll want to know: \hi{खंभा}}
\textbf{\hi{खंभा}} — \emph{khambhā} — \textquotedblleft{}a pole; a pillar\textquotedblright{}.

\begin{quote}
\hi{बिजली} \hi{का} \hi{खंभा} \hi{सड़क} \hi{पर} \hi{गिर} \hi{गया।}
\end{quote}

\begin{quote}
\emph{bijlī kā khambhā saṛak par gir gayā.}
\end{quote}

\begin{quote}
The electricity pole fell onto the road.
\end{quote}

Masculine.

\begin{cousinweb}
\textbf{\hi{खंभा}} (\emph{khambhā}) is a tadbhava of Sanskrit \emph{skambha}, \textquotedblleft{}a prop, a pillar\textquotedblright{} (CDIAL 13639), which lists Hindi \emph{khambh}, \textquotedblleft{}a pillar, a pole\textquotedblright{}; sk- became kh-.
\end{cousinweb}

\begin{grammarlens}[title={What you've built}]
One more word for getting around.
\end{grammarlens}

\subsection*{Your turn}
\begin{itemize}
\raggedright
  \item \emph{Say it:} \emph{khambhā}
  \item \emph{Say it:} \emph{khambhā}, once more
  \item \emph{Say it:} \emph{bijlī kā khambhā saṛak par gir gayā.} — The electricity pole fell onto the road.
  \item \emph{Recall:} say the Hindi for a thatched roof, then the Hindi for a guard, then say \emph{khambhā} again
\end{itemize}

\begin{tcolorbox}[breakable,colback=teal!4,colframe=teal!35!black,title={Before you move on}]
What does \hi{खंभा} mean? (\textquotedblleft{}A pole; a pillar\textquotedblright{}.) Say it once more.
\end{tcolorbox}
\section[\texorpdfstring{\hi{शोर} (śor)}{शोर (śor)}]{\hi{शोर} (śor) — noise}

\label{lesson:HI-C423-sor}



\begin{quote}
Before the new one: say the Hindi for a guard, then the Hindi for a pole.
\end{quote}

\subsection*{You'll want to know: \hi{शोर}}
\textbf{\hi{शोर}} — \emph{śor} — \textquotedblleft{}noise\textquotedblright{}.

\begin{quote}
\hi{सड़क} \hi{पर} \hi{बहुत} \hi{शोर} \hi{है।}
\end{quote}

\begin{quote}
\emph{saṛak par bahut śor hai.}
\end{quote}

\begin{quote}
It is very noisy on the road.
\end{quote}

Masculine.

\begin{cousinweb}
\textbf{\hi{शोर}} (\emph{śor}) is usually traced to Persian \emph{šor}, \textquotedblleft{}noise, clamour\textquotedblright{}.
\end{cousinweb}

\begin{grammarlens}[title={What you've built}]
One more word for getting around.
\end{grammarlens}

\subsection*{Your turn}
\begin{itemize}
\raggedright
  \item \emph{Say it:} \emph{śor}
  \item \emph{Say it:} \emph{śor}, once more
  \item \emph{Say it:} \emph{saṛak par bahut śor hai.} — It is very noisy on the road.
  \item \emph{Recall:} say the Hindi for a guard, then the Hindi for a pole, then say \emph{śor} again
\end{itemize}

\begin{tcolorbox}[breakable,colback=teal!4,colframe=teal!35!black,title={Before you move on}]
What does \hi{शोर} mean? (\textquotedblleft{}Noise\textquotedblright{}.) Say it once more.
\end{tcolorbox}
\section[\texorpdfstring{\hi{गुर्राना} (gurrānā)}{गुर्राना (gurrānā)}]{\hi{गुर्राना} (gurrānā) — to growl}

\label{lesson:HI-C423-gurrana}



\begin{quote}
Before the new one: say the Hindi for a pole, then the Hindi for noise.
\end{quote}

\subsection*{You'll want to know: \hi{गुर्राना}}
\textbf{\hi{गुर्राना}} — \emph{gurrānā} — \textquotedblleft{}to growl\textquotedblright{}.

\begin{quote}
\hi{रात} \hi{को} \hi{गली} \hi{का} \hi{कुत्ता} \hi{गुर्रा} \hi{रहा} \hi{था।}
\end{quote}

\begin{quote}
\emph{rāt ko galī kā kuttā gurrā rahā thā.}
\end{quote}

\begin{quote}
At night the dog in the lane was growling.
\end{quote}

Past (it …): \hi{वह} \hi{शेर} \hi{गुर्राया}, \emph{vah śer gurrāyā}.

\begin{cousinweb}
\textbf{\hi{गुर्राना}} (\emph{gurrānā}) goes back to \emph{*guragura}, \textquotedblleft{}growl\textquotedblright{}, a form CDIAL reconstructs and treats as imitative (CDIAL 4207), which lists Hindi \emph{gurrānā}: the word rumbles like the sound.
\end{cousinweb}

\begin{grammarlens}[title={What you've built}]
One more word for getting around.
\end{grammarlens}

\subsection*{Your turn}
\begin{itemize}
\raggedright
  \item \emph{Say it:} \emph{gurrānā}
  \item \emph{Say it:} \emph{gurrānā}, once more
  \item \emph{Say it:} \emph{rāt ko galī kā kuttā gurrā rahā thā.} — At night the dog in the lane was growling.
  \item \emph{Recall:} say the Hindi for a pole, then the Hindi for noise, then say \emph{gurrānā} again
\end{itemize}

\begin{tcolorbox}[breakable,colback=teal!4,colframe=teal!35!black,title={Before you move on}]
What does \hi{गुर्राना} mean? (\textquotedblleft{}To growl\textquotedblright{}.) Say it once more.
\end{tcolorbox}
\section[\texorpdfstring{\hi{भौंकना} (bhaũknā)}{भौंकना (bhaũknā)}]{\hi{भौंकना} (bhaũknā) — to bark}

\label{lesson:HI-C423-bhaukna}



\begin{quote}
Before the new one: say the Hindi for noise, then the Hindi for to growl.
\end{quote}

\subsection*{You'll want to know: \hi{भौंकना}}
\textbf{\hi{भौंकना}} — \emph{bhaũknā} — \textquotedblleft{}to bark\textquotedblright{}.

\begin{quote}
\hi{रात} \hi{को} \hi{गली} \hi{के} \hi{कुत्ते} \hi{बहुत} \hi{भौंकते} \hi{हैं।}
\end{quote}

\begin{quote}
\emph{rāt ko galī ke kutte bahut bhaũkte haiṁ.}
\end{quote}

\begin{quote}
The dogs in the lane bark a lot at night.
\end{quote}

Past (it …): \hi{वह} \hi{कुत्ता} \hi{रात} \hi{भर} \hi{भौंका}, \emph{vah kuttā rāt bhar bhaũkā}.

\begin{cousinweb}
\textbf{\hi{भौंकना}} (\emph{bhaũknā}) is a tadbhava of Sanskrit \emph{bukkati}, \textquotedblleft{}barks\textquotedblright{} (CDIAL 9265), which lists Hindi \emph{bhaũknā}.
\end{cousinweb}

\begin{grammarlens}[title={What you've built}]
One more word for getting around.
\end{grammarlens}

\subsection*{Your turn}
\begin{itemize}
\raggedright
  \item \emph{Say it:} \emph{bhaũknā}
  \item \emph{Say it:} \emph{bhaũknā}, once more
  \item \emph{Say it:} \emph{rāt ko galī ke kutte bahut bhaũkte haiṁ.} — The dogs in the lane bark a lot at night.
  \item \emph{Recall:} say the Hindi for noise, then the Hindi for to growl, then say \emph{bhaũknā} again
\end{itemize}

\begin{tcolorbox}[breakable,colback=teal!4,colframe=teal!35!black,title={Before you move on}]
What does \hi{भौंकना} mean? (\textquotedblleft{}To bark\textquotedblright{}.) Say it once more.
\end{tcolorbox}
